\section{Soundness proofs}
\label{app:proofs}

This appendix proves that each structural check is sound: if the check passes, the
corresponding safety property holds for all valid execution traces.

\begin{lemma}[Exit Reachability Soundness]
\label{lem:exit_reach}
If $\mathrm{ExitReach}(G)$ holds (i.e., the BFS/DFS from $v_0$ visits
all of $V_f$), then for every exit node $v_f \in V_f$, there exists a
directed path from $v_0$ to $v_f$ in $G$.
\end{lemma}
\begin{proof}
The check computes $\mathrm{Reach}(v_0)$ by BFS from $v_0$, which
correctly identifies all nodes reachable via directed edges
\citep{clarke1999modelchecking}. If $v_f \in \mathrm{Reach}(v_0)$,
then by the correctness of BFS there exists a path $v_0 \to \cdots \to
v_f$. Since the check passes only when $V_f \subseteq
\mathrm{Reach}(v_0)$, the result follows.
\end{proof}

\begin{lemma}[Dead-End Soundness]
\label{lem:dead_end}
If $\mathrm{NoDead}(G)$ holds, then every non-exit node in every trace
of $G$ has at least one successor, and hence no maximal trace terminates
at a non-exit node.
\end{lemma}
\begin{proof}
Let $v$ be any node with $\kappa_V(v) \neq \textsc{exit}$. Since
$\mathrm{NoDead}(G)$ holds, there exists $u$ such that $(v,u) \in E$.
Therefore, if a trace $\pi$ visits $v$, the trace can be extended by
$u$. A trace can only be maximal (i.e., cannot be extended) when it
ends at a node in $V_f$ or at a node with no outgoing edges. Since
every non-exit node has outgoing edges, all maximal traces end at exit
nodes.
\end{proof}

\begin{lemma}[Router Shape Soundness]
\label{lem:router_shape}
If $\mathrm{RouterShape}(G)$ holds, then every outgoing transition from
a router node is explicitly labeled as conditional.
\end{lemma}
\begin{proof}
Direct from the predicate definition: the check iterates all edges
$(v,u) \in E$ where $\kappa_V(v) = \textsc{router}$ and verifies
$\kappa_E(v,u) = \textsc{conditional}$. Passing the check establishes
the universal quantification.
\end{proof}

\begin{lemma}[Human Gate Soundness]
\label{lem:human_gate}
If $\mathrm{HumanGate}(G)$ holds, then the graph contains at least one
node classified as a human-in-the-loop step.
\end{lemma}
\begin{proof}
The check searches for $v \in V$ with $\kappa_V(v) = \textsc{human}$.
Passing the check certifies existence.
\end{proof}

\begin{lemma}[Tool Declaration Soundness]
\label{lem:tool_decl}
If $\mathrm{ToolDecl}(G)$ holds, then every tool node explicitly
declares its tool set.
\end{lemma}
\begin{proof}
For each $v$ with $\kappa_V(v) = \textsc{tool}$, the check verifies
$T(v) \neq \emptyset$. Passing the check establishes the universal
statement.
\end{proof}

\begin{lemma}[Reverse Reachability Soundness]
\label{lem:reverse_reach}
If $\mathrm{ExitReachAll}(G)$ holds (i.e.,
$\mathrm{Reach}(v_0) \subseteq \mathrm{RevReach}(V_f)$), then for every
node~$v$ reachable from~$v_0$, there exists a directed path from~$v$ to
some exit node~$v_f \in V_f$.
\end{lemma}
\begin{proof}
The check computes $\mathrm{RevReach}(V_f)$ by BFS on the reverse
adjacency starting from all nodes in~$V_f$. By the correctness of BFS,
$v \in \mathrm{RevReach}(V_f)$ implies the existence of a path
$v \to \cdots \to v_f$ for some $v_f \in V_f$. The check passes only when
$\mathrm{Reach}(v_0) \subseteq \mathrm{RevReach}(V_f)$, so the result
holds for every reachable node.
\end{proof}

\begin{lemma}[Human Gate Coverage Soundness]
\label{lem:human_gate_cov}
If $\mathrm{HumanGateCov}(G, S)$ holds for a set~$S$ of sensitive tool
names, then every path from~$v_0$ to a node invoking a tool in~$S$
passes through at least one \textsc{human} node.
\end{lemma}
\begin{proof}
The check constructs a modified graph~$G'$ by removing all
\textsc{human}-typed nodes and their incident edges, then computes
$\mathrm{Reach}_{G'}(v_0)$.  If no sensitive tool node is in
$\mathrm{Reach}_{G'}(v_0)$, then in the original graph~$G$ every path
from~$v_0$ to such a node must traverse a removed \textsc{human} node.
Conversely, if a sensitive tool node is reachable in~$G'$, the
check fails and reports the human-free path.
\end{proof}

\begin{lemma}[Static Temporal Soundness]
\label{lem:static_temporal}
For a compiled monitor rule with DFA $(Q, q_0, \delta, F_{\mathrm{viol}})$
and an agent workflow graph~$G$, if the graph~$\times$~DFA product
construction reports no violation, then no execution trace of~$G$
violates the temporal property.
\end{lemma}
\begin{proof}
The product BFS explores all reachable states in
$V \times Q$.  Each product state $(v, q)$ represents being at graph
node~$v$ with the DFA in state~$q$.  For each successor~$u$ of~$v$ in~$G$,
the DFA transitions via $q' = \delta(q, \sigma(v))$ where $\sigma(v)$ is the
event symbol for node~$v$, and the product successor $(u, q')$ is enqueued.
If no reachable product state has $q' \in F_{\mathrm{viol}}$, then no
execution trace --- which corresponds to a path through product states ---
can reach a violation state.  This is an over-approximation: the product
explores all \emph{graph paths}, which is a superset of actual runtime traces
(since runtime traces depend on LLM decisions at router nodes).  Hence the
analysis is sound but potentially conservative.
\end{proof}

\begin{lemma}[Temporal Monitor Soundness]
\label{lem:monitor}
For each of the three base DFA patterns (forbidden, implication-future,
until), if the compiled monitor reports no violation on a finite trace
$\pi$, then $\pi$ satisfies the corresponding temporal formula.
\end{lemma}
\begin{proof}
The proof proceeds by cases.

\emph{Forbidden ($\mathbf{G}\,\neg a$):} The DFA has two states:
$s_0$ (safe) and $s_1$ (violated, absorbing). The transition
function maps $s_0 \xrightarrow{a} s_1$ and
$s_0 \xrightarrow{\neg a} s_0$, with $s_1$ absorbing.
If the monitor never enters $s_1$, then $a$ was false at every step,
establishing $\mathbf{G}\,\neg a$.

\emph{Implication-future ($a \to \mathbf{F}\,b$):} The DFA has three
states: $s_0$ (idle), $s_1$ (waiting for $b$), and $s_2$ (violated,
absorbing). Transition: $s_0 \xrightarrow{a \wedge \neg b} s_1$;
$s_1 \xrightarrow{b} s_0$; $s_1 \xrightarrow{a \wedge \neg b} s_2$.
If the monitor never enters $s_2$, then every occurrence of $a$ is
followed by $b$ before $a$ recurs.

\emph{Until ($a\,\mathbf{U}\,b$):} The DFA has three states: $s_0$
(waiting), $s_1$ (satisfied, absorbing), $s_2$ (violated, absorbing).
Transition: $s_0 \xrightarrow{b} s_1$;
$s_0 \xrightarrow{a \wedge \neg b} s_0$;
$s_0 \xrightarrow{\neg a \wedge \neg b} s_2$.
If the monitor never enters $s_2$, then $a$ held on every step until
$b$ occurred, establishing $a\,\mathbf{U}\,b$.

In each case the DFA state invariant is maintained by induction on
the trace prefix.
\end{proof}
